\documentclass[11pt,a4paper]{article}
\usepackage[T1]{fontenc}
\usepackage[utf8]{inputenc}
\usepackage{lmodern,amsmath,amssymb,amsthm,booktabs,tabularx}
\usepackage[margin=27mm]{geometry}
\usepackage[hidelinks]{hyperref}
\hypersetup{pdftitle={Classical spins and the dynamical meaning of an operational state},pdfauthor={Author-review draft}}
\setlength{\emergencystretch}{3em}
\newtheorem{proposition}{Proposition}
\newcommand{\E}{\mathbb E}
\newcommand{\R}{\mathbb R}
\title{Classical spins and the dynamical meaning of an operational state}
\author{Author-review draft}
\date{13 September 2026}
\begin{document}
\maketitle
\begin{abstract}
A classical random orientation observed through affine directional responses
has the operational state space of a qubit. We use two interacting classical
spins to show why this finite statistical description requires a separate
dynamical consistency condition. Two preparations with identical retained
statistics acquire different terminal probabilities under one smooth reversible
Hamiltonian. The interaction admits no exact finite enlargement of the retained
expectation coordinates for all probability preparations, even with nonlinear
coordinate updates. The explicit example distinguishes statistical dimension,
distinguishability capacity and reversible operational closure, and identifies
the additional physical input required when applying quantum reconstruction
principles to classical mechanics.
\end{abstract}

\section{When a state description predicts its own evolution}
An operational state groups preparations that give the same probabilities for
an assigned collection of transformations and terminal measurements. Enlarging
that collection can split the equivalence classes \cite{preparation}.
Here we start with a quotient defined by terminal tests and local rotations,
then ask whether a classical interaction preserves it. A state update on this
quotient exists precisely when equal retained states remain equal after the
interaction. This tests the sufficiency of the description independently of
microscopic reversibility.

For two spins the retained data are their mean orientations and their nine
cross moments: $\E u_i$, $\E v_j$ and $\E[u_i v_j]$. These fifteen numbers
predict every initial product test in the model. We exhibit two probability
preparations with the same fifteen numbers but different later measurement
probabilities, then prove that adding finitely many expectation coordinates
cannot fix this for all preparations. This is the concrete contribution of
the worked example.

Section 2 defines the measurement model and its relation to qubit statistics;
Section 3 gives the separating preparations and their explicit motion.
Section 4 proves the finite-enlargement obstruction. The counterexample itself
needs only classical spin evolution and expectations; the final obstruction
also uses finite-dimensional linear algebra.

Restricted classical models and minimal tensor products provide a standard
way to separate local statistics from composition \cite{janotta}. The local
sphere-to-ball reduction is an instance of the classical-extension construction
reviewed by Bacciagaluppi \cite{bacciagaluppi}. Quantum
reconstruction results make composition and reversible transformations
substantive hypotheses \cite{hardy,cdp,torre}. We illustrate their physical
role with a classical spin interaction and an exact observable-closure test.
Finite invariant observable spaces are familiar from Koopman theory
\cite{brunton}; the argument here is an explicit application to the retained
measurement statistics of this model. Unlike the nonlinear quantum boxes of
Ref.~\cite{preparation}, our full classical probability measures evolve by
linear pushforward; the failure occurs when projecting onto retained moments.

The example uses all Borel probability preparations and stipulated affine
terminal measurements. Its result concerns exact finite expectation-coordinate
descriptions on a nonzero time interval. Approximate closure and restricted
preparations pose different questions. The fixed spin magnitudes supply action
units; the construction leaves the physical selection of a universal action
constant as a separate task.

\section{A classical orientation with qubit statistics}
Let the hidden orientation be $u\in S^2\subset\R^3$. A preparation is any
Borel probability measure $\mu$. An allowed measurement is a finite family
\begin{equation}
 e_i(u)=a_i+b_i\cdot u,\qquad a_i\ge |b_i|,\qquad
 \sum_i a_i=1,\quad\sum_i b_i=0.
\end{equation}
These nonnegative responses sum to one. Their probabilities depend exactly
on $r=\int u\,d\mu$, since directional binary measurements determine its
three components. Every $|r|\le1$ is a mixture of opposite orientations.
The state space is therefore the unit ball, with three normalized coordinates
and four unnormalized linear coordinates.

Using Pauli matrices as a coordinate basis,
\begin{equation}
 \rho(r)=\tfrac12(I+r\cdot\sigma),\qquad E_i=a_iI+b_i\cdot\sigma,
 \qquad \operatorname{tr}(E_i\rho)=a_i+b_i\cdot r.
\end{equation}
Thus the allowed local states and effects have the usual qubit representation.
Rotations $u\mapsto Ru$, $R\in SO(3)$, descend to $r\mapsto Rr$ and
connect all pure ball states continuously and reversibly.

The perfectly distinguishable capacity is two. If outcomes distinguish $k$
states $r_i$ perfectly, then $1=a_i+b_i\cdot r_i\le2a_i$ for each success
outcome. Summing gives $k\le2$; opposite poles attain two. Capacity counts
perfectly distinguishable alternatives in one measurement, whereas statistical
dimension counts independent probability coordinates.

\subsection{Composition supplies another choice}
For $n$ components admit all joint measures on $(S^2)^n$ and product
measurements, their mixtures and coarse-grainings. The retained state is
\begin{equation}
 z=\int (1,u_1)\otimes\cdots\otimes(1,u_n)\,d\mu.
\end{equation}
Product effects span these $4^n$ unnormalized coordinates. The state space is
the convex hull of product ball states, equivalently separable density
matrices in the Pauli representation. Local rotations and permutations
preserve this description.

Its capacity is $2^n$: every outcome is bounded by $2^n$ times its value at
the product of zero ball vectors, including after mixing and coarse-graining.
Normalization bounds the number of perfectly distinguished states by $2^n$;
product pole tests attain it. Every extreme composite state is product-pure,
because the compact product-pure set generates its convex hull. Conversely,
pure marginals force each orientation almost surely, making such products
extreme. Consequently every marginal of a pure composite is pure, so a mixed
local state has no purification in this composite class.

There is also a simple subspace contrast. For two components the equal-sign
effect along axis $n$ is
\begin{equation}
 e_{\rm same}(u,v)=\tfrac12[1+(n\cdot u)(n\cdot v)].
\end{equation}
Probability one forces support on $(n,n)$ and $(-n,-n)$. The resulting
capacity-two face is a segment rather than an elementary ball. This illustrates
the additional content of Hardy's subspace requirement \cite{hardy} and of
purification in the CDP reconstruction \cite{cdp}.

\section{A reversible interaction distinguishes equal retained states}
Give the spins fixed magnitudes $S_A,S_B>0$ with action units and choose
$J>0$ with energy units. Standard classical-spin Poisson brackets
\cite{spins} give
\begin{equation}
 \{u_i,u_j\}=S_A^{-1}\epsilon_{ijk}u_k,\quad
 \{v_i,v_j\}=S_B^{-1}\epsilon_{ijk}v_k,\quad\{u_i,v_j\}=0,
 \qquad H=Ju_zv_z.
\end{equation}
Use $\dot f=\{f,H\}$ and $\Omega_A=J/S_A$, $\Omega_B=J/S_B$.
The equations include
\begin{equation}\label{eq:flow}
 \dot u_x=-\Omega_Au_yv_z,\quad\dot u_y=\Omega_Au_xv_z,
 \quad\dot u_z=0,
\end{equation}
with the corresponding equations for $v$ obtained by exchanging the spins.
Both sphere norms and $H$ are conserved. Smoothness on the compact product
of spheres gives a complete flow $\Phi_t$ with inverse $\Phi_{-t}$.

Let $e_x,e_y,e_z$ be coordinate unit vectors and prepare
\begin{align}
 \mu_A&=\delta_{e_y}(du)\,\tfrac12(\delta_{e_z}+\delta_{-e_z})(dv),\\
 \mu_B&=\delta_{e_y}(du)\,\tfrac12(\delta_{e_x}+\delta_{-e_x})(dv).
\end{align}
Both have $\E u=e_y$, $\E v=0$ and $\E uv^T=0$, so every retained
product measurement gives identical initial statistics. Both preparations
also have microscopic energy zero.

Since $u_z=0$, the second spin remains fixed. In $\mu_A$, writing
$v=s e_z$ with $s=\pm1$,
\begin{equation}
 u(t)=(-s\sin\Omega_A t,\cos\Omega_A t,0).
\end{equation}
In $\mu_B$ we have $v_z=0$, so both orientations remain fixed. The allowed
product effect $e(u,v)=(1+u_x)(1+v_z)/4$ therefore gives
\begin{equation}\label{eq:separation}
 p_A(t)=\tfrac14(1-\sin\Omega_A t),\qquad p_B(t)=\tfrac14.
\end{equation}
For sufficiently small positive time the outputs differ despite identical
retained inputs. Thus even a single-valued reduced-state update fails for
this interaction and preparation class. The hidden flow remains reversible,
and each evolved preparation still defines a valid separable operational state.
The missing information appears directly in
\begin{equation}
 \frac{d}{dt}\E[u_xv_z]=-\Omega_A\E[u_yv_z^2].
\end{equation}

\section{Exact finite enlargement also fails}
The question is whether storing more averages could make the reduced state
sufficient. For exact prediction, every evolved retained observable must have
an expectation determined by those stored averages. The first proposition
exhibits infinitely many independent evolved observables; the second explains
why nonlinear processing of finitely many averages cannot recover them for
every preparation.

The pullback $U_tf=f\circ\Phi_t$ acts on bounded Borel observables compared
pointwise. All atomic preparations are admitted, so this is stronger than
equality modulo a chosen reference measure.

\begin{proposition}
No finite-dimensional real space of bounded Borel observables containing
$1$ and the original terminal effects is invariant under $U_t$ for every
time in any nonzero interval, including a forward interval beginning at zero.
\end{proposition}
\begin{proof}
The complexified space contains $w=u_x+iu_y$. Equation~\eqref{eq:flow} and
constancy of $v_z$ give $U_tw=w\exp(i\Omega_A t v_z)$. Suppose finitely
many distinct times obey $\sum_{j=1}^m c_jU_{t_j}w=0$. Restrict to
$u=e_x$ and $v=(\sqrt{1-z^2},0,z)$, $-1<z<1$. Then
\begin{equation}
 \sum_{j=1}^m c_j\exp(i\Omega_A t_jz)=0.
\end{equation}
Differentiating in $z$ at zero through order $m-1$ gives a Vandermonde system
in the distinct numbers $i\Omega_A t_j$. Its determinant is nonzero since
$\Omega_A>0$, so every coefficient vanishes. There are arbitrarily many
independent time translates, contradicting finite dimension. Only explicit
exponentials were differentiated; added observables need no regularity.
\end{proof}

\begin{proposition}
Allowing nonlinear updates of a finite expectation vector does not repair
exact evolution for all admitted preparations.
\end{proposition}
\begin{proof}
Let $f_0=1,f_1,\ldots,f_d$ be retained real observables. If $\E_\mu g$ is
determined by their expectations for every probability measure $\mu$, then
$g$ lies in their real linear span. Otherwise evaluation linear algebra gives
finitely many points $x_j$ and real weights $a_j$ such that
\begin{equation}
 \sum_j a_j f_k(x_j)=0\quad(0\le k\le d),\qquad
 \sum_j a_jg(x_j)\ne0.
\end{equation}
Indeed, if every relation between the evaluation vectors also annihilated
$g$, evaluation of $g$ would define a linear functional on their span and
express $g$ as a combination of the $f_k$. The relation for $f_0=1$ gives
equal positive and negative masses. Normalizing them produces two atomic
probability measures with identical retained expectations but different
$g$-expectations. Apply this observation to $g=U_tf_k$: any single-valued
expectation update forces invariance of the retained span, which the preceding
proposition excludes.
\end{proof}

These propositions concern finite statistical dimension with ordinary
probabilistic mixing. Infinite observable dimension by itself says nothing
about perfectly distinguishable capacity. Approximation, restrictions on
preparations and changed dynamics each alter a stated premise. At $J=0$
the interaction disappears and static finite closure is restored.

\section{What the example contributes to reconstruction}
De la Torre and colleagues \cite{torre} study locally quantum, locally
tomographic systems with continuous reversible operational interactions.
Their theorem constrains composite structure and, with the specified
identical-copy and operational closure assumptions, selects quantum theory.
The spin example already has the required local ball representation. Its
Hamiltonian interaction instead fails to induce a map on that finite
operational state, and finite enlargement cannot supply one for all the
admitted preparations. Microscopic reversibility therefore does not discharge
the theorem's operational reversibility hypothesis.

\begin{center}\small
\begin{tabularx}{\linewidth}{@{}lX@{}}\toprule
Premise & Role in this example\\\midrule
Affine terminal effects & Stipulated restriction yielding the local ball.\\
All Borel preparations & Makes equality of expectations a pointwise test and permits atomic separation.\\
Product measurements & Selects the minimal composite and its finite moment coordinates.\\
Classical spin Hamiltonian & Supplies a complete reversible microscopic evolution.\\
Exact finite operational closure & Fails under that interaction for the stated preparation class.\\
Physical action identification & Spin magnitudes are supplied; universality requires additional physical input.\\\bottomrule
\end{tabularx}
\end{center}

The dimensionless ball construction fixes probability geometry. Mechanical
time enters the interaction through $\Omega_A t=Jt/S_A$, where $J$ and
$S_A$ are independently specified model parameters. Relating a reconstructed
transformation generator to measured energy and a universal phase per unit
action is a further physical identification. The useful conclusion is a
precise location for that requirement, alongside the requirement of a finite
description stable under the allowed dynamics.

\begin{thebibliography}{9}
\bibitem{preparation} E. G. Cavalcanti, N. C. Menicucci and J. L. Pienaar,
``The preparation problem in nonlinear extensions of quantum theory,''
\href{https://arxiv.org/abs/1206.2725v1}{arXiv:1206.2725v1} (2012), pp. 2--4.
\bibitem{bacciagaluppi} G. Bacciagaluppi, ``Classical Extensions, Classical
Representations and Bayesian Updating in Quantum Mechanics,''
\href{https://arxiv.org/abs/quant-ph/0403055v1}{arXiv:quant-ph/0403055v1}
(2004), Sec. 2.3.
\bibitem{janotta} P. Janotta and H. Hinrichsen, ``Generalized Probability
Theories: What determines the structure of quantum theory?''
\href{https://arxiv.org/abs/1402.6562v3}{arXiv:1402.6562v3} (2014),
Secs. 4 and 5.1.
\bibitem{hardy} L. Hardy, ``Quantum Theory From Five Reasonable Axioms,''
\href{https://arxiv.org/abs/quant-ph/0101012v4}{arXiv:quant-ph/0101012v4}
(2001), Secs. 1, 4 and 7.
\bibitem{cdp} G. Chiribella, G. M. D'Ariano and P. Perinotti,
``Informational Derivation of Quantum Theory,'' Phys. Rev. A \textbf{84},
012311 (2011), \href{https://arxiv.org/abs/1011.6451v3}{arXiv:1011.6451v3}.
\bibitem{torre} G. de la Torre, Ll. Masanes, A. J. Short and M. P. M\"uller,
``Deriving quantum theory from its local structure and reversibility,''
\href{https://arxiv.org/abs/1110.5482v1}{arXiv:1110.5482v1} (2011),
Theorems 1--2.
\bibitem{spins} S. Rado\v{s}evi\'c et al., ``Hamiltonian Dynamics of Classical
Spins,'' \href{https://arxiv.org/abs/2503.16308v2}{arXiv:2503.16308v2}
(2025), Secs. V.1--V.2.
\bibitem{brunton} S. L. Brunton, B. W. Brunton, J. L. Proctor and J. N. Kutz,
``Koopman Invariant Subspaces and Finite Linear Representations of Nonlinear
Dynamical Systems for Control,'' PLOS ONE \textbf{11}(2), e0150171 (2016),
\href{https://doi.org/10.1371/journal.pone.0150171}{doi:10.1371/journal.pone.0150171}.
\end{thebibliography}
\end{document}
